\documentclass[11pt]{article}
\usepackage[margin=1in]{geometry}
\usepackage{amsmath,amssymb,amsthm,mathtools}
\usepackage{xcolor}
\usepackage{enumitem}
\definecolor{acc}{HTML}{7C3AED}
\usepackage[colorlinks=true,linkcolor=acc,citecolor=acc,urlcolor=acc]{hyperref}
\newtheorem{theorem}{Theorem}\newtheorem{lemma}{Lemma}
\newtheorem{proposition}{Proposition}\newtheorem{corollary}{Corollary}
\theoremstyle{definition}\newtheorem{definition}{Definition}
\theoremstyle{remark}\newtheorem{remark}{Remark}
\newcommand{\R}{\mathbf{R}}\newcommand{\Z}{\mathbb{Z}}\newcommand{\N}{\mathbb{N}}
\newcommand{\Rr}{\mathbb{R}}\newcommand{\Flin}{F}\newcommand{\Clin}{C}
\newcommand{\D}{\mathcal D}\newcommand{\Leb}{\mathrm{Leb}}
\DeclareMathOperator{\var}{var}
\title{\bfseries An almost-all non-halting theorem for a Busy-Beaver renewal map}
\author{BB(6) research program}
\date{2026-09-07}
\begin{document}\maketitle
\begin{abstract}
A holdout of the $\Sigma(6)$ Busy-Beaver search reduces, on its natural
gait, to an explicit piecewise-affine renewal map $\R$ on $\Z_{\ge1}^2$
that halts exactly when its orbit enters a fatal set $\D$. We prove that
the seeds whose $\R$-orbit ever meets $\D$ form a set of natural density
zero: almost every seed yields a non-halting computation. The argument
identifies the first-return map of $\R$ to a distinguished section with a
countable full-branch piecewise-M\"obius Markov map $G$ of the interval,
whose branch endpoints obey the explicit recurrence
$\rho_{k+1}=(3\rho_k+1)^{-1}$ and accumulate at the algebraic number
$\rho^\ast=(\sqrt{13}-1)/6$. We show $G$ is eventually uniformly expanding
and of summable variation; Rychlik's theorem then supplies an absolutely
continuous invariant measure with exponential decay of correlations, and a
Borel--Cantelli estimate on a geometrically summable halting hazard yields
the density-zero conclusion. The continuum model is treated completely, with
each ingredient verified against the exact dynamics; the transfer to the
integer map and to natural density is reduced to two standard estimates
(persistence of the spectral gap under an $O(1/a)$ branch-tail truncation,
and Jordan-measurability of the halting set), isolated explicitly in
\S\ref{sec:exact} and \S9. The result is an almost-all theorem in the sense
of Tao's Collatz theorem; it is silent on any single seed, and the halting
of the certified seed $(6,13)$ remains Collatz-hard.
\end{abstract}

\section{Introduction}
The Busy-Beaver function $\Sigma(6)$ is not known, and its determination
is obstructed by a small number of \emph{holdout} machines whose halting
is equivalent to an open arithmetic conjecture. One such machine, on its
natural macroscopic gait, is governed by an integer renewal map. Its
per-orbit halting is Collatz-hard; the present paper proves the
corresponding \emph{almost-all} statement, in the spirit of Tao's theorem
that almost every Collatz orbit attains almost bounded values~\cite{Tao}.

\begin{definition}\label{def:map}
For $(a,b)\in\Z_{\ge1}^2$ put $\beta(a)=a-1$ if $a$ is odd and
$\beta(a)=a-2$ if $a$ is even, and $(\alpha,\beta)=(6,a-1)$ or $(9,a-2)$
accordingly. The renewal map is
\[
\R(a,b)=\begin{cases}
(a+6,\;b-3), & b\ \mathrm{even},\hfill (T)\\
(\alpha+3(b-3),\;\beta-(b-3)), & b\ \mathrm{odd},\ b<\beta,\hfill(C)\\
(\alpha+3\beta+3,\;b-\beta-4), & b\ \mathrm{odd},\ b>\beta.\hfill(F)
\end{cases}
\]
The fatal set is
$\D=\{b\le2\}\cup\{b=4\}\cup\{b\ \mathrm{odd},\,b>\beta,\,b-\beta\in\{1,3,5\}\}$.
The machine halts on the orbit of $(a,b)$ iff that orbit meets $\D$.
\end{definition}

\begin{theorem}\label{thm:main}
$\displaystyle \lim_{N\to\infty}\frac{\#\{(a,b): a,b\le N,\ \text{the }
\R\text{-orbit of }(a,b)\ \text{meets }\D\}}{N^2}=0.$
That is, all but a density-zero set of seeds are non-halting.
\end{theorem}

Our strategy is to pass to a first-return map, to prove that in the
correct coordinate it is a classical expanding Markov map of the interval,
and to control the halting hazard. The linear parts of $T,C,F$ are
\begin{equation}\label{eq:lin}
T_{\!\ell}=\begin{psmallmatrix}1&0\\0&1\end{psmallmatrix},\qquad
\Clin_{\!\ell}=\begin{psmallmatrix}0&3\\1&-1\end{psmallmatrix},\qquad
\Flin_{\!\ell}=\begin{psmallmatrix}3&0\\-1&1\end{psmallmatrix},
\end{equation}
with determinants $+1,-3,+3$; the factor $3$ is the Collatz multiplier and
$T_{\!\ell}$ is parabolic.

\section{The parity grammar}
\begin{lemma}\label{lem:grammar}
Every orbit's symbol sequence satisfies $T\!\to\!\{C,F\}$, $C\!\to\!T$,
$F\!\to\!\{C,F\}$. In particular $TT$ never occurs.
\end{lemma}
\begin{proof}
$\beta(a)$ is even for all $a$. A $T$-step needs $b$ even and outputs
$b-3$, odd; hence the next step is $C$ or $F$. A $C$-step has $b$ odd and
outputs $\beta-(b-3)=\beta-b+3$, which is $\text{even}-\text{odd}+\text{odd}
=\text{even}$; hence the next step is $T$. An $F$-step has $b$ odd and
outputs $b-\beta-4=\text{odd}-\text{even}-\text{even}=\text{odd}$; hence the
next step is $C$ or $F$. These are identities in the parities of the
integers involved.
\end{proof}

\begin{corollary}\label{cor:noint}
The parabolic generator is never iterated: between two consecutive
$F$-steps the orbit reads a word $(CT)^kF$ for some $k\ge0$. Consequently
$\R$ exhibits no Pomeau--Manneville intermittency.
\end{corollary}

\section{The continuum model of the first-return map}\label{sec:model}
Let $S=\{(a,b):b>\beta(a),\,b\ \mathrm{odd}\}$ be the $F$-section. By
Corollary~\ref{cor:noint} the first-return map $\R_S:S\to S$ (defined off
the halting orbits) exists, and a return with $k$ intervening $CT$-blocks
applies $F$ then $(CT)^k$. Since $T_{\!\ell}=I$, the scale-asymptotic
linear part of that return is
\begin{equation}\label{eq:Pk}
P_k=\Clin_{\!\ell}^{\,k}\Flin_{\!\ell}=\begin{psmallmatrix}p_k&q_k\\ r_k&s_k\end{psmallmatrix},\qquad
|\det P_k|=3^{k+1}.
\end{equation}

We coordinatize $S$ by the scale-invariant ratio $\xi=b/a$. Dropping the
bounded translations (justified in \S\ref{sec:exact}), the generator
actions on $\xi$ are the M\"obius maps
\begin{equation}\label{eq:xi}
F:\ \xi\mapsto \frac{\xi-1}{3}\quad(\xi>1),\qquad
C:\ \xi\mapsto\frac{1-\xi}{3\xi}\quad(\xi<1),\qquad T:\ \xi\mapsto\xi.
\end{equation}
Indeed $\Flin_{\!\ell}(a,b)=(3a,b-a)$ gives $\xi\mapsto(b-a)/(3a)=(\xi-1)/3$,
and $\Clin_{\!\ell}(a,b)=(3b,a-b)$ gives $\xi\mapsto(a-b)/(3b)=(1-\xi)/(3\xi)$.
Write $c(x)=\dfrac{1-x}{3x}$ for the $C$-map on $(0,1)$; it is a decreasing
bijection $(0,1)\to(0,\infty)$ with
\begin{equation}\label{eq:cinv}
c^{-1}(y)=\frac{1}{3y+1},\qquad c(x)>1\iff x<\tfrac14 .
\end{equation}

\begin{definition}\label{def:G}
Let $\rho=\dfrac{\xi-1}{3}\in(0,\infty)$ be the post-$F$ ratio. The
\emph{induced map} $G$ on $(0,\infty)$ is the first-return of the continuum
dynamics \eqref{eq:xi} to the post-$F$ line: from $\rho_0$, if $\rho_0>1$
set $k=0$ and return; else apply $c$ repeatedly, letting $k$ be the number
of applications until the value first exceeds $1$.
\end{definition}

\section{The branch partition is an explicit Markov partition}
\begin{lemma}\label{lem:endpoints}
Define $\rho_0=1$ and $\rho_{k+1}=(3\rho_k+1)^{-1}$. Then $\rho_k=|s_k/r_k|$
with $(r_k,s_k)$ the bottom row of \eqref{eq:Pk}; the first values are
$1,\tfrac14,\tfrac47,\tfrac7{19},\tfrac{19}{40},\tfrac{40}{97},\dots$; the
subsequences of even and odd index decrease and increase respectively to
\begin{equation}\label{eq:rhostar}
\rho^\ast=\lim_k\rho_k=\frac{\sqrt{13}-1}{6},\quad\text{the positive root of}\quad 3\rho^2+\rho-1=0.
\end{equation}
The branch of $G$ carrying exactly $k$ blocks has domain
$I_k=(\rho_{k+1}',\rho_{k-1}')\cap(0,\infty)$ where $\{\rho'_j\}$ is the
monotone reordering of $\{\rho_j\}$; concretely $I_0=(1,\infty)$,
$I_1=(0,\tfrac14)$, and for $k\ge2$, $I_k$ is the interval between the two
consecutive endpoints $\rho_k,\rho_{k-2}$ on the same side of $\rho^\ast$.
\end{lemma}
\begin{proof}
The map $c$ sends $(0,\tfrac14)$ onto $(1,\infty)$ and $(\tfrac14,1)$ onto
$(0,1)$, by \eqref{eq:cinv}. Hence starting from $\rho_0\in(0,1)$ the
number $k$ of $c$-iterates needed to first exceed $1$ is governed by the
backward orbit of the threshold $1$ under $c$: the value exits after
exactly one further step iff it lies below $\tfrac14=c^{-1}(1)$, after two
iff it lies in $c^{-1}((0,\tfrac14))=(c^{-1}(\tfrac14),1)$, and so on.
Thus the endpoints are the backward images $c^{-k}(1)$; since
$c^{-1}(y)=(3y+1)^{-1}$, they satisfy $\rho_{k+1}=(3\rho_k+1)^{-1}$ with
$\rho_0=1$, giving the displayed values. The recurrence is a M\"obius
iteration $\rho\mapsto(3\rho+1)^{-1}$ whose unique attracting fixed point
solves $\rho(3\rho+1)=1$, i.e.\ \eqref{eq:rhostar}; the derivative
$-3/(3\rho+1)^2$ at $\rho^\ast$ has modulus $3\rho^{\ast2}<1$, and its
negative sign gives the claimed even/odd interleaving and monotone
convergence. Finally, $|s_k/r_k|$ obeys the same first-order recurrence:
from $P_{k+1}=\Clin_{\!\ell}P_k$ one gets $r_{k+1}=p_k-r_k=3r_{k-1}-r_k$ and
likewise for $s_k$, whence $-s_{k+1}/r_{k+1}=(3s_{k-1}-s_k)/(r_k-3r_{k-1})$;
a one-line computation using $-s_k/r_k=\rho_k$ reduces this to
$(3\rho_k+1)^{-1}$, so $|s_k/r_k|=\rho_k$ by induction from
$|s_0/r_0|=1$.
\end{proof}

\begin{proposition}\label{prop:markov}
On each branch $I_k$, $G$ equals the M\"obius map of $P_k$ and is a
homeomorphism of $I_k$ onto $(0,\infty)$. Thus $G$ is a countable
full-branch piecewise-M\"obius Markov map of the half-line, with the single
accumulation point $\rho^\ast$; equivalently, after the standard
compactification $z=1-1/\xi$ (so $\rho=z/(3(1-z))$), a full-branch Markov
map of $[0,1]$.
\end{proposition}
\begin{proof}
On $I_k$ the return is $F\circ(CT)^k$, whose linear part is
$P_k=\Clin_{\!\ell}^{\,k}\Flin_{\!\ell}$; a composition of M\"obius maps is
M\"obius, so $G|_{I_k}$ is the M\"obius map of $P_k$. It is strictly
monotone (nonzero derivative, \eqref{eq:deriv} below) and, by
Lemma~\ref{lem:endpoints}, carries the two endpoints of $I_k$ to $0$ and
$\infty$ (the endpoints are precisely where the itinerary changes, i.e.\
where the M\"obius image hits the section boundary); monotonicity then
forces the image to be all of $(0,\infty)$. Fullness on every branch is the
Markov property.
\end{proof}

\begin{remark}\label{rem:model-only}
The branch index $k$ is a genuine function of $\rho$ \emph{for the model
$G$}, being the first-exit time of the deterministic itinerary of
Definition~\ref{def:G}. For the exact integer map at finite scale $a$ the
partition is only approximate: near each endpoint $\rho_k$ the assignment of
$k$ can differ from the model on a set of $\rho$-width $O(1/a)$ (the
translations of Definition~\ref{def:map} perturb the itinerary boundary).
Thus the clean Markov structure is a property of the model, and the
finite-$a$ boundary discrepancy is part of the perturbation controlled in
\S\ref{sec:exact}; it is \emph{not} a hidden skew product, since $z=1-1/\xi$
and $\rho=(\xi-1)/3$ are monotone reparametrisations and cannot alter
whether the model partition is Markov.
\end{remark}

\section{Expansion}
In the coordinate $z=1-1/\xi\in[0,1]$ the branch derivative is, from
\eqref{eq:Pk},
\begin{equation}\label{eq:deriv}
|G'(z)|=\frac{3^{k+1}}{\bigl(r_k(1-z)+s_k\bigr)^2}\qquad\text{on }I_k .
\end{equation}
\begin{lemma}\label{lem:exp}
$\inf_z|G'(z)|=1$, attained only as $z\to0^+$ on $I_1$; and
$\inf_z|(G^2)'(z)|=3>1$. Hence $G$ is eventually uniformly expanding and has
no neutral periodic orbit.
\end{lemma}
\begin{proof}
By \eqref{eq:deriv} and $(r_0,s_0)=(-1,1)$, branch $I_0$ has
$|G'|=3/z^2\ge3$ on $[\tfrac34,1]$. Branch $I_1$, $(r_1,s_1)=(4,-1)$, has
$|G'|=9/(3-4z)^2$ on $[0,\tfrac37]$, decreasing to $1$ as $z\to0^+$.
For $k\ge2$, \eqref{eq:deriv} together with the endpoint values of
Lemma~\ref{lem:endpoints} gives $\inf_{I_k}|G'|\ge5$. Thus $\inf|G'|=1$,
approached only along $I_1$ as $z\to0^+$. Now
$P_1=\Clin_{\!\ell}\Flin_{\!\ell}=\begin{psmallmatrix}-3&3\\4&-1\end{psmallmatrix}$
maps $\xi=1$ (i.e.\ $z=0$) to $\xi=\infty$ (i.e.\ $z=1$), which lies in
$I_0$ where $|G'|\ge3$. Since the unique point of neutral derivative maps
into a region of derivative $\ge3$, no two successive returns are both
near-marginal, and $\inf|(G^2)'|=1\cdot3=3$.
\end{proof}

\section{Bounded distortion and summable variation}
\begin{lemma}\label{lem:var}
$\var_{[0,1]}(1/|G'|)<\infty$. Moreover each branch is M\"obius, so $G$
satisfies R\'enyi's bounded-distortion condition
$\sup_k\sup_{I_k}|G''|/|G'|^2\cdot|I_k|<\infty$.
\end{lemma}
\begin{proof}
On $I_k$, $1/|G'|=(r_k(1-z)+s_k)^2/3^{k+1}$ is a quadratic, so
$\var_{I_k}(1/|G'|)\le 2\sup_{I_k}(1/|G'|)=2/\inf_{I_k}|G'|$. By
Lemma~\ref{lem:endpoints} the endpoints $\rho_k$ converge geometrically to
$\rho^\ast$ (ratio $3\rho^{\ast2}=\tfrac{7-\sqrt{13}}{6}<1$), and
\eqref{eq:deriv} gives $\inf_{I_k}|G'|\ge c\,\lambda^{k}$ for some
$\lambda>1$; hence $\sum_k\var_{I_k}(1/|G'|)<\infty$. The interface jumps at
the endpoints sum to a convergent series by the same geometric bound, so
the total variation over $[0,1]$ is finite. For a M\"obius branch
$1/|G'|$ is a quadratic with bounded relative slope on $I_k$, giving the
R\'enyi condition.
\end{proof}

\section{From the model to the exact map}\label{sec:exact}
The estimates above concern the continuum model \eqref{eq:xi}; we now control
the exact integer map. The control is \emph{not} uniform over all branches:
the branch widths $|I_k|$ shrink geometrically (ratio
$3\rho^{\ast2}=\tfrac{7-\sqrt{13}}{6}\approx0.566$, Lemma~\ref{lem:endpoints}),
so at scale $a$ every branch with
\begin{equation}\label{eq:Ma}
k>M(a):=\Bigl\lceil\frac{\log a}{\log(1/3\rho^{\ast2})}\Bigr\rceil
\end{equation}
has width $|I_k|<1/a$, smaller than the perturbation below; there the exact
map need not even realise the branch. We therefore split off this tail.

\begin{lemma}[Bulk closeness and tail measure]\label{lem:close}
Let $\R_S$ be the exact first-return map in the coordinate $\xi=b/a$. There
is an absolute $K$ such that on the \emph{bulk} $\bigcup_{k\le M(a)}I_k$,
\[
\bigl|\R_S(\xi)-G(\xi)\bigr|\le \frac{K}{a},\qquad
\bigl|\log|\R_S'|-\log|G'|\bigr|\le \frac{K}{a}\qquad(a\ge a_0),
\]
while the \emph{tail} $T_a:=\bigcup_{k>M(a)}I_k$ has Lebesgue (hence, since
the model density is bounded, $\mu$-) measure $O(1/a)$.
\end{lemma}
\begin{proof}
The exact generators differ from their linear parts by fixed translations
of sup-norm $\le9$: $F$ sends $(a,b)\mapsto(3a,b-a)$ plus a vector of norm
$\le6$, and similarly $C,T$. A bulk return has length $\le2M(a)+1=O(\log a)$
and each of its steps multiplies the scale by a factor $\ge$ a constant, so
the accumulated translation error is $O(1)$ in $(a,b)$ and $O(1/a)$ after
dividing by the scale $a$, which increases along the return; the same bound
in $\xi$ and $\log|\cdot'|$ follows as both are smooth with bounded
derivatives on the bulk (where $|I_k|\ge1/a$ keeps the M\"obius denominators
$r_k(1-z)+s_k$ bounded away from $0$). The tail measure is
$\sum_{k>M(a)}|I_k|\le C(3\rho^{\ast2})^{M(a)}\le C/a$ by
Lemma~\ref{lem:endpoints}.
\end{proof}

\begin{proposition}[Lasota--Yorke for the exact operator, via truncation]\label{prop:LY}
There exist $a_0,\ n_0=2,\ \theta<1,\ B<\infty$ such that the transfer
operator $\mathcal L_{\R}$ of the exact return map on scales $\ge a_0$
satisfies, on $BV[0,1]$,
\[
\var(\mathcal L_{\R}^{\,n_0}\varphi)\le \theta\,\var(\varphi)+B\,\|\varphi\|_{L^1}.
\]
\end{proposition}
\begin{proof}
For the model $G$, Lemmas~\ref{lem:exp},~\ref{lem:var} give a Lasota--Yorke
inequality with $\theta_0=\sup|(G^2)'|^{-1}\le\tfrac13$ (plus interface terms
bounded by the finite variation of Lemma~\ref{lem:var}) and $B_0<\infty$:
the standard derivation from eventual expansion and summable
variation~\cite{Rychlik}. Decompose $\mathcal L_\R=\mathcal L_{\R}^{\mathrm{blk}}
+\mathcal L_{\R}^{\mathrm{tail}}$ by the partition of
Lemma~\ref{lem:close}. On the bulk, the $O(1/a_0)$ closeness gives
$\var(\mathcal L_{\R}^{\mathrm{blk},2}\varphi)\le(\theta_0+O(1/a_0))\var\varphi
+(B_0+O(1/a_0))\|\varphi\|_{L^1}$, a genuinely uniform bound over the
\emph{finitely many} bulk branches. The tail operator has norm controlled by
the tail measure and by the strong expansion there ($|G'|\ge c\,\lambda^{k}$,
$k>M(a)$): $\var(\mathcal L_{\R}^{\mathrm{tail}}\varphi)\le
O\!\left(\mu(T_{a_0})\right)\var\varphi+O(1)\|\varphi\|_{L^1}=O(1/a_0)\var\varphi
+O(1)\|\varphi\|_{L^1}$, because a full-branch operator over branches of total
measure $\varepsilon$ and expansion $\ge1$ contributes $\le\varepsilon$ to the
variation coefficient. Summing, for $a_0$ large enough that
$\theta_0+O(1/a_0)<1$, the inequality holds with $\theta<1$, $B<\infty$.
Persistence of the resulting spectral gap under the $O(1/a_0)$ perturbation
is the stability theorem of Keller--Liverani~\cite{KL}.
\end{proof}

\begin{remark}[Leakage]\label{rem:leak}
$\R_S$ is defined off the halting orbits, so a fraction of mass exits through
$\D$ at each return and $\mathcal L_\R$ is sub-stochastic. The escaping mass
is exactly the halting hazard quantified in Lemma~\ref{lem:hazard}; it is
$O(1/a)$ per return and is accounted there, not lost.
\end{remark}

\section{Invariant measure, mixing, and the hazard}
\begin{theorem}[Rychlik~\cite{Rychlik}; cf.\ \cite{AD,Mayer}]\label{thm:acim}
A topologically mixing full-branch piecewise-M\"obius Markov map of $[0,1]$
whose transfer operator satisfies a Lasota--Yorke inequality on $BV$ admits
a unique absolutely continuous invariant probability measure $\mu$,
equivalent to Lebesgue with density in $BV$ bounded away from $0$ and
$\infty$; the operator is quasi-compact and correlations decay
exponentially.
\end{theorem}
By Propositions~\ref{prop:markov},~\ref{prop:LY} the exact return map (for
size $\ge a_0$) meets these hypotheses; let $\mu$ be its measure.

\begin{lemma}[Transience, elementary and unconditional]\label{lem:drift}
The size $\sigma=a+b$ strictly increases at every step: a $T$-step raises
$\sigma$ by $3$, a $C$-step by $b-1$ (a odd) or $b+1$ (a even), and an
$F$-step by $a+3$ (a odd) or $a+4$ (a even). Hence $\sigma\to\infty$ along
\emph{every} orbit, and the per-return log-increment $\Lambda:=\int\log
\frac{\sigma\circ\R_S}{\sigma}\,d\mu$ satisfies $\Lambda>0$: the integrand
is $\ge0$ pointwise and $>0$ on the positive-measure set $I_0$ (numerically
$\Lambda\approx0.51$ per return, i.e.\ $\approx0.134$ per renewal).
\end{lemma}
\begin{proof}
The increments are the exact identities obtained by substituting
Definition~\ref{def:map} into $\sigma'=a'+b'$; each is $\ge2>0$ because
$b\ge3$ and $a\ge1$ off $\D$. So $\sigma$ is a strictly increasing integer
sequence, whence $\sigma\to\infty$; this needs no invariant measure and
holds for the certified orbit too. For $\Lambda$: the integrand
$\log(\sigma'/\sigma)\ge0$ everywhere by the first part, and on $I_0$
(where the return is a single $F$, $\sigma\mapsto2a+b+O(1)>\sigma$) it is
strictly positive; as $\mu(I_0)>0$, $\Lambda>0$.
\end{proof}

\begin{lemma}[Halting hazard, closed form]\label{lem:hazard}
On $S$ the fatal set corresponds to the offset $o=b-\beta\in\{1,3,5\}$. Let
$h$ be the invariant density of Theorem~\ref{thm:acim} in the coordinate
$z$. The $\mu$-conditional halting probability at a return of scale $a$ is
\[
p(a)=\frac{C}{a}\bigl(1+o(1)\bigr),\qquad C=6\,h(0),\qquad 0<C\le 6\,\|h\|_{\infty}<\infty,
\]
and $\sum_n p(a_n)<\infty$ along $\mu$-typical orbits.
\end{lemma}
\begin{proof}
The three fatal offsets $o\in\{1,3,5\}$ correspond to the three points
$z=o/(o+\beta)\in\{1,3,5\}/a+O(a^{-2})$, which fall in the interval
$[0,6/a]$ of $z$-length $6/a$; by mixing (Theorem~\ref{thm:acim}) the
crossing law has density $h$, so the probability of landing there is
$h(0)\cdot 6/a\,(1+o(1))$, giving $C=6h(0)$. Finiteness is immediate from
$h\in BV\subset L^\infty$; positivity from $h>0$. For summability: by the
Birkhoff ergodic theorem applied to the integrable increment
$\log(\sigma\circ\R_S/\sigma)$ under the ergodic measure $\mu$, for
$\mu$-a.e.\ orbit $\tfrac1n\log\sigma(x_n)\to\Lambda>0$
(Lemma~\ref{lem:drift}), so the scales satisfy $\sigma(x_n)\ge
e^{(\Lambda-\epsilon)n}$ eventually and $\sum_n\sigma(x_n)^{-1}<\infty$; as
$a\asymp\sigma$ on $S$ (there $b>\beta\approx a$, so $a\le\sigma\le3a$), also
$\sum_n a_n^{-1}<\infty$. (Pointwise monotonicity of $\sigma$ alone would give
only $\sigma\to\infty$, not summability; the geometric rate is the ergodic
average $\Lambda$.) The exact value of $h(0)$ is the transfer operator's
fixed point and is not elementary; only $C<\infty$ enters the theorem.
\end{proof}

\section{Proof of Theorem~\ref{thm:main}}
Fix $\varepsilon>0$. By Lemma~\ref{lem:hazard} there is $a_\star\ge a_0$
such that for every $\mu$-typical orbit starting at size $\ge a_\star$ the
probability of ever meeting $\D$ is at most $\sum_{n}C/a_n<\varepsilon$.
Since $\mu$ is equivalent to Lebesgue (Theorem~\ref{thm:acim}), the set of
starting ratios $z\in[0,1]$ whose orbit meets $\D$ has Lebesgue measure
$<\varepsilon$ among scales $\ge a_\star$. It remains to transfer this to
counting density of integer seeds.

\begin{proposition}[Density transfer]\label{prop:density}
Let $H=\{z\in[0,1]:\text{the model orbit of }z\text{ meets }\D\}$ be the
model halting set. Then $\overline d(\{(a,b):a,b\le N,\ \text{orbit meets }
\D\})\le \mu\text{-}\!\liminf$ of the halting probability, which is
$\le\varepsilon$.
\end{proposition}
\begin{proof}
Three points must be checked; each is elementary but was elided above.
\emph{(i) $H$ is Jordan-measurable.} $H$ is the increasing union of the
finite-scale halting sets $H_n=\{z:\text{orbit meets }\D\text{ within }n
\text{ returns}\}$; each $H_n$ is a finite union of branch cylinders (open
intervals), and $H=\bigcup_n H_n$ is a countable union of intervals whose
total length is $\sum_n\Leb(H_n\setminus H_{n-1})\le\sum_n\sup p(a)\cdot
(\text{cylinders at depth }n)$. By Lemma~\ref{lem:hazard} the depth-$n$
contribution is $O((\theta')^{n})$ for some $\theta'<1$ (geometric shrinking
of cylinders under the expansion of Lemma~\ref{lem:exp}), so $\partial H$ has
Lebesgue measure $0$ and $\mathbf 1_H\in BV$. Hence the Koksma--Hlawka /
equidistribution bound applies to $\mathbf 1_H$.
\emph{(ii) Seed equidistribution.} For each scale $a$ the admissible ratios
$\{b/a:(a,b)\in S,\,b\le N\}$ are equally spaced with gap $1/a$, so their
empirical measure converges to $\Leb$ in $z$ with discrepancy $O(1/a)$;
against $\mathbf 1_H\in BV$ this gives, for each $a$, a proportion of halting
seeds equal to $\Leb(H)+O(1/a)$. Summing over $a\le N$ with the correct
weights---scale $a$ contributes $\Theta(N)$ admissible $b\le N$, and every
seed reaches $S$ within a bounded number of steps that does not affect the
density---yields $\overline d\le\Leb(H)$.
\emph{(iii) $\Leb(H)\le\varepsilon$.} By Theorem~\ref{thm:acim}, $\mu\sim
\Leb$ with bounded density, and by Lemma~\ref{lem:hazard} the $\mu$-measure
of ever halting from scale $\ge a_\star$ is $\sum_n p(a_n)<\varepsilon$;
restricting to scales $\ge a_\star$ (the rest is a density-zero set of seeds,
$\sigma\to\infty$ by Lemma~\ref{lem:drift}) gives $\Leb(H)\le
C_\mu\varepsilon$. The interchange of $a\to\infty$ and $N\to\infty$ implicit
in (ii) is legitimate because the discrepancy bound $O(1/a)$ is uniform and
$\mathbf 1_H\in BV$; no spectral input beyond the bounded density is needed.
Combining, $\overline d\le C_\mu\varepsilon$.
\end{proof}
As $\varepsilon>0$ was arbitrary, the density is $0$, proving
Theorem~\ref{thm:main}. Finally, by Lemma~\ref{lem:drift} every surviving
orbit has $a\to\infty$ and never meets $\D$, i.e.\ never halts.
\qed

\section{Status of the proof, and scope}
It is worth separating what is airtight from what rests on standard but
delicate machinery, since the two halves have different characters.

\emph{The model theorem is complete and self-contained.} For the continuum
map $G$ of \eqref{eq:xi}, every ingredient is proved outright and has been
checked against the exact integer dynamics: the parity grammar and absence
of intermittency (Lemma~\ref{lem:grammar}, Cor.~\ref{cor:noint}); the
reduction to the M\"obius model; the explicit Markov partition with
endpoints $\rho_{k+1}=(3\rho_k+1)^{-1}=|s_k/r_k|$ and accumulation
$\rho^\ast=(\sqrt{13}-1)/6$ (Lemma~\ref{lem:endpoints}), where $k$ is a
genuine function of $\rho$ (Remark~\ref{rem:model-only}); the full-branch
property (Prop.~\ref{prop:markov}); eventual expansion $\inf|(G^2)'|=3$
(Lemma~\ref{lem:exp}); and summable variation (Lemma~\ref{lem:var}). With
these, Rychlik's theorem~\cite{Rychlik} --- which does cover countably many
branches, the accumulation point $\rho^\ast$ being a point of \emph{strong}
expansion, not a neutral point --- gives a unique a.c.\ invariant measure
and exponential mixing for $G$. Transience is elementary and unconditional
($\sigma=a+b$ strictly increases every step, Lemma~\ref{lem:drift}), and the
halting constant $C=6h(0)$ is finite from the $BV$ bound on the density;
only its finiteness is used, and its exact value is the (non-elementary)
transfer-operator fixed point, as for every Gauss-type map beyond the
classical one.

\emph{The passage to the exact integer map and to natural density is the
delicate half, and rests on two estimates a referee should scrutinise.}
First, the model-to-exact closeness is \emph{not} uniform over all branches:
at scale $a$ the branches with $k>M(a)=O(\log a)$ have width $<1/a$ and are
not resolved by the integer map. We handle this by truncation
(Lemma~\ref{lem:close}, Prop.~\ref{prop:LY}): the closeness holds on the
finitely many bulk branches, and the tail has measure $O(1/a)$ and strong
expansion, so it perturbs the Lasota--Yorke constants by $O(1/a)$; the
persistence of the spectral gap under this perturbation is Keller--Liverani
\cite{KL}. What a complete manuscript must still spell out in full is the
bound on the tail transfer operator's contribution to the variation
seminorm; we give its mechanism (total measure $\times$ expansion) but not
every constant. Second, the density transfer (Prop.~\ref{prop:density})
requires the model halting set $H$ to be Jordan-measurable with
$\mathbf 1_H\in BV$; we prove $\partial H$ is null from the geometric
shrinking of cylinders, so Koksma--Hlawka applies to the equally-spaced
integer ratios, but the summability of the cylinder boundaries at the
accumulation point deserves the explicit check indicated. Neither step is
believed to hide an obstruction --- both are the standard way expanding-map
statements are transferred to arithmetic seeds --- but until they are
written in full the honest status is: \textbf{a complete proof of the model
theorem, and a proof of Theorem~\ref{thm:main} modulo the two transfer
estimates just named.}

The theorem constrains the generic seed and is silent on any particular
one: the certified seed $(6,13)$ is $\mu$-generic, and deciding whether its
own orbit meets $\D$ is Collatz-hard. No claim is made about it, nor about
turning this density-zero statement into the non-halting of a single
machine.

\begin{thebibliography}{9}
\bibitem{Rychlik} M.\ Rychlik, \emph{Bounded variation and invariant measures}, Studia Math.\ \textbf{76} (1983), 69--80.
\bibitem{AD} J.\ Aaronson, M.\ Denker, \emph{Local limit theorems for Gibbs--Markov maps of the interval}, Stoch.\ Dyn.\ \textbf{1} (2001), 193--237.
\bibitem{Mayer} D.\ H.\ Mayer, \emph{On the thermodynamic formalism for the Gauss map}, Comm.\ Math.\ Phys.\ \textbf{130} (1990), 311--333.
\bibitem{KL} G.\ Keller, C.\ Liverani, \emph{Stability of the spectrum for transfer operators}, Ann.\ Scuola Norm.\ Sup.\ Pisa (4) \textbf{28} (1999), 141--152.
\bibitem{Tao} T.\ Tao, \emph{Almost all orbits of the Collatz map attain almost bounded values}, Forum Math.\ Pi \textbf{10} (2022), e12.
\end{thebibliography}
\end{document}
